\subsection*{The Extended Addition Operation}

\begin{definitionbox}{Definition (Addition on \texorpdfstring{$\mathbb{W}$}{W})}
\begin{definition}[Addition on $\mathbb{W}$]
\label{def:addition-on-w}
Define $+_{\mathbb{W}}:\mathbb{W}\times\mathbb{W}\to\mathbb{W}$ by
\[
0+_{\mathbb{W}}a=a,\qquad
a+_{\mathbb{W}}0=a,
\]
and, for $a,b\in\mathbb{N}$,
\[
a+_{\mathbb{W}}b=a+_{\mathbb{N}}b.
\]
\end{definition}
\end{definitionbox}

\begin{remark*}[Standard quantified statement]
\[
\begin{aligned}
&+_{\mathbb{W}}:\mathbb{W}\times\mathbb{W}\to\mathbb{W},\\
&\forall a\in\mathbb{W}\;(0+_{\mathbb{W}}a=a\ \wedge\ a+_{\mathbb{W}}0=a),\\
&\forall a,b\in\mathbb{N}\;
  (a+_{\mathbb{W}}b=a+_{\mathbb{N}}b).
\end{aligned}
\]
\end{remark*}

\begin{remark*}[Interpretation]
If either input is $0$, the new identity clause applies. If both inputs lie in
$\mathbb{N}$, the old operation is used. Since $0\notin\mathbb{N}$, the
positive-positive clause does not overlap with the zero clauses.
\end{remark*}

\begin{dependencies}
\begin{itemize}
  \item \hyperref[def:whole-numbers]{Whole Numbers}
  \item \hyperref[def:addition-on-peano-system]{Addition on a Peano System}
\end{itemize}
\end{dependencies}

\begin{theorembox}{Theorem (Addition on \texorpdfstring{$\mathbb{W}$}{W} Is Well-Defined)}
\begin{theorem}[Addition on $\mathbb{W}$ Is Well-Defined]
\label{thm:addition-on-w-well-defined}
The case definition of $+_{\mathbb{W}}$ determines a binary operation on
$\mathbb{W}$.
\hyperref[prf:addition-on-w-well-defined]{Go to proof.}
\end{theorem}
\end{theorembox}

\begin{remark*}[Standard quantified statement]
\[
+_{\mathbb{W}}\text{ is a well-defined binary operation on }\mathbb{W}.
\]
\end{remark*}

\begin{remark*}[Interpretation]
Well-definedness is a case check. The zero cases are new, the
positive-positive case is inherited from $\mathbb{N}$, and compatibility is
automatic because the positive-positive case cannot contain $0$.
\end{remark*}

\begin{dependencies}
\begin{itemize}
  \item \hyperref[def:whole-numbers]{Whole Numbers}
  \item \hyperref[def:addition-on-w]{Addition on $\mathbb{W}$}
  \item \hyperref[thm:addition-well-defined-on-peano-system]{Addition Is Well-Defined on a Peano System}
\end{itemize}
\end{dependencies}

\begin{theorembox}{Theorem (Zero Is an Additive Identity on \texorpdfstring{$\mathbb{W}$}{W})}
\begin{theorem}[Zero Is an Additive Identity on $\mathbb{W}$]
\label{thm:zero-additive-identity-on-w}
For every $a\in\mathbb{W}$,
\[
0+_{\mathbb{W}}a=a
\qquad\text{and}\qquad
a+_{\mathbb{W}}0=a.
\]
\hyperref[prf:zero-additive-identity-on-w]{Go to proof.}
\end{theorem}
\end{theorembox}

\begin{remark*}[Standard quantified statement]
\[
\forall a\in\mathbb{W}\;
\bigl(0+_{\mathbb{W}}a=a\ \wedge\ a+_{\mathbb{W}}0=a\bigr).
\]
\end{remark*}

\begin{remark*}[Interpretation]
This is the structural feature missing from the one-based natural numbers.
After adjoining $0$, addition has a genuine identity element.
\end{remark*}

\begin{dependencies}
\begin{itemize}
  \item \hyperref[def:addition-on-w]{Addition on $\mathbb{W}$}
\end{itemize}
\end{dependencies}

\begin{theorembox}{Theorem (Addition on \texorpdfstring{$\mathbb{W}$}{W} Extends Addition on \texorpdfstring{$\mathbb{N}$}{N})}
\begin{theorem}[Addition on $\mathbb{W}$ Extends Addition on $\mathbb{N}$]
\label{thm:addition-on-w-extends-addition-on-n}
For all $a,b\in\mathbb{N}$,
\[
a+_{\mathbb{W}}b=a+_{\mathbb{N}}b.
\]
\hyperref[prf:addition-on-w-extends-addition-on-n]{Go to proof.}
\end{theorem}
\end{theorembox}

\begin{remark*}[Standard quantified statement]
\[
\forall a,b\in\mathbb{N}\;
\bigl(a+_{\mathbb{W}}b=a+_{\mathbb{N}}b\bigr).
\]
\end{remark*}

\begin{remark*}[Interpretation]
The extension does not alter sums of positive natural numbers.
\end{remark*}

\begin{dependencies}
\begin{itemize}
  \item \hyperref[def:addition-on-w]{Addition on $\mathbb{W}$}
  \item \hyperref[def:addition-on-peano-system]{Addition on a Peano System}
\end{itemize}
\end{dependencies}

\begin{theorembox}{Theorem (Addition on \texorpdfstring{$\mathbb{W}$}{W} Is Associative)}
\begin{theorem}[Addition on $\mathbb{W}$ Is Associative]
\label{thm:addition-on-w-associative}
For all $a,b,c\in\mathbb{W}$,
\[
(a+_{\mathbb{W}}b)+_{\mathbb{W}}c
=
a+_{\mathbb{W}}(b+_{\mathbb{W}}c).
\]
\hyperref[prf:addition-on-w-associative]{Go to proof.}
\end{theorem}
\end{theorembox}

\begin{remark*}[Standard quantified statement]
\[
\forall a,b,c\in\mathbb{W}\;
\bigl((a+_{\mathbb{W}}b)+_{\mathbb{W}}c
=a+_{\mathbb{W}}(b+_{\mathbb{W}}c)\bigr).
\]
\end{remark*}

\begin{remark*}[Interpretation]
Associativity follows by reducing zero cases to the identity law and reducing
the positive-positive-positive case to associativity on $\mathbb{N}$.
\end{remark*}

\begin{dependencies}
\begin{itemize}
  \item \hyperref[thm:zero-additive-identity-on-w]{Zero Is an Additive Identity on $\mathbb{W}$}
  \item \hyperref[thm:addition-associative]{Addition Is Associative}
\end{itemize}
\end{dependencies}

\begin{theorembox}{Theorem (Addition on \texorpdfstring{$\mathbb{W}$}{W} Is Commutative)}
\begin{theorem}[Addition on $\mathbb{W}$ Is Commutative]
\label{thm:addition-on-w-commutative}
For all $a,b\in\mathbb{W}$,
\[
a+_{\mathbb{W}}b=b+_{\mathbb{W}}a.
\]
\hyperref[prf:addition-on-w-commutative]{Go to proof.}
\end{theorem}
\end{theorembox}

\begin{remark*}[Standard quantified statement]
\[
\forall a,b\in\mathbb{W}\;
\bigl(a+_{\mathbb{W}}b=b+_{\mathbb{W}}a\bigr).
\]
\end{remark*}

\begin{remark*}[Interpretation]
The zero cases are symmetric by definition, and the positive-positive case is
the commutativity theorem already proved for $\mathbb{N}$.
\end{remark*}

\begin{dependencies}
\begin{itemize}
  \item \hyperref[def:addition-on-w]{Addition on $\mathbb{W}$}
  \item \hyperref[thm:addition-commutative]{Addition Is Commutative}
\end{itemize}
\end{dependencies}

\begin{theorembox}{Theorem (Addition on \texorpdfstring{$\mathbb{W}$}{W} Satisfies Cancellation)}
\begin{theorem}[Addition on $\mathbb{W}$ Satisfies Cancellation]
\label{thm:addition-on-w-cancellation}
For all $a,b,c\in\mathbb{W}$,
\[
a+_{\mathbb{W}}b=a+_{\mathbb{W}}c\Longrightarrow b=c.
\]
\hyperref[prf:addition-on-w-cancellation]{Go to proof.}
\end{theorem}
\end{theorembox}

\begin{remark*}[Standard quantified statement]
\[
\forall a,b,c\in\mathbb{W}\;
\bigl(a+_{\mathbb{W}}b=a+_{\mathbb{W}}c\Longrightarrow b=c\bigr).
\]
\end{remark*}

\begin{remark*}[Interpretation]
When $a=0$, cancellation is the identity law. When all relevant entries are
positive, it is the natural-number cancellation theorem.
\end{remark*}

\begin{dependencies}
\begin{itemize}
  \item \hyperref[thm:zero-additive-identity-on-w]{Zero Is an Additive Identity on $\mathbb{W}$}
  \item \hyperref[thm:addition-left-cancellation]{Left Cancellation for Addition}
\end{itemize}
\end{dependencies}
